\documentclass[12pt]{book}

\usepackage[dvips,letterpaper,margin=0.75in,bottom=0.5in]{geometry}
\usepackage{cite}
\usepackage{slashed}
\usepackage{graphicx}
\usepackage{mathtools}
\usepackage{amsmath}
\usepackage{amssymb}
\usepackage{braket}

\DeclareMathOperator{\sinc}{sinc}

\begin{document}


\title{PHY 115L \\
Introduction to the \\
Discrete Fourier Transform \\ 
\author{Michael Mulhearn}}

\maketitle

\setcounter{chapter}{1}
\chapter{Discrete Fourier Transform}

\section{Fourier Series for Complex-Valued Functions}

We will treat the Fourier Series in the context of the inner product space of square-integrable periodic functions of period $a$,  with an inner product defined as:
$$ \braket{f|g} \coloneqq \int_{-\frac{a}{2}}^{\frac{a}{2}} f^*(x) g(x) dx$$
The complex exponential functions
$$e_k(x) \coloneqq \sqrt{\frac{1}{a}} \exp\left(i \frac{2\pi k x}{a}\right)$$
are orthonormal vectors.  To verify:
\begin{eqnarray*}
  \braket{e_j|e_k} &=& \frac{1}{a} \int_{-\frac{a}{2}}^{\frac{a}{2}}
  \exp\left(i \frac{2\pi (k-j) x}{a}\right) dx \\[7pt]
  &=& \frac{1}{\pi (k-j)} \cdot \frac{\exp\left(i \pi (k-j) \right)-\exp\left(-i \pi (k-j) \right)}{2i}\\[7pt]
  &=& \sinc(\pi (k-j))
\end{eqnarray*}
As $k$ and $j$ are integers, and $\sinc$ is zero for every integer multiple of $\pi$ except for $sin(0)=1$, we have shown:
\begin{equation}
  \braket{e_j|e_k} = \delta_{kj}
\end{equation}
We will not derive it here, but the complex functions are also a complete basis, so we can expand any function in terms of the basis of complex exponentials by:
\begin{equation}
 \psi(x) = \sum_{j=-\infty}^{+\infty} c_j e_j (x)
\end{equation}
Because of othonormality, we can determine the Fourier coeeficients by:
\begin{equation}
c_k = \braket{e_k | \psi}
\end{equation}

\section{Discrete Fourier Transform}

Next we want to consider the case where we only know the value of $\psi(x)$ at $n$ discrete positions of $x$, specifically:
$$x_j = x_0 + \lambda j$$
for $x_0=-a/2$, $j=0,1,\ldots,n-1$ and $\lambda$ the interval between data points:
$$\lambda = \frac{a}{n}$$
Notice that the sequence ends before $x=a/2$ because this point is redundant to $x=-a/2$ for the periodic functions that we are considering.

We account for the discrete sampling by making the replacement:
$$\psi(x) \to \sum_{m=0}^{n-1} \delta(x-x_j) \lambda f(x)$$
which is choosen because:
$$\int_{x_1}^{x_2} \sum_{m=0}^{n-1} \delta(x-x_j) \lambda f(x) dx = \sum_{m=m_1}^{m_2} \lambda f_m \; \xrightarrow[\lambda \to 0]{} \; \int_{x_1}^{x_2} f(x) dx $$
showing that integration (smoothing) of discrete sampled function
reproduces the integral of the original function as the space between
samples approaches zero.  We do expect this sampling to introduce
high-frequency artifacts which must be considered, which we will
discuss later.

Calculating the fourier series coefficients with the sampled function we obtain:
\begin{eqnarray*}
  c_k &=& \braket{e_k | \psi} \\[5pt]
  &=& \sqrt{\frac{1}{a}} \int_{-\frac{a}{2}}^{\frac{a}{2}}
  \exp\left(-i \frac{2\pi k x}{a}\right) \sum_{m=0}^{n-1} \delta(x-x_j) \lambda f(x) dx \\[7pt]
  &=& \sqrt{\frac{\lambda}{n}} \sum_{m=0}^{n-1}
  \exp\left(-i \frac{2\pi m k}{n}\right) \exp\left(-i \frac{2\pi k x_0}{a}\right) f_j dx \\[7pt]
\end{eqnarray*}
which we can reorganize as:
\begin{equation}
  \frac{c_k}{\sqrt{\lambda}} \exp\left(i \frac{2\pi k x_0}{a}\right)
 = \frac{1}{\sqrt{n}} \sum_{m=0}^{n-1}
  \exp\left(-i \frac{2\pi m k}{n}\right) f_m \equiv A_k  
\end{equation}
where $A_k$ are the Fourier coefficients from {\tt numpy.fft.fft} if called with the normalization option {\tt norm=``ortho''}.  The complex exponential on the left hand side is a phase factor that accounts the location of the first data point $x_0$.  Notice that the right hand side depends only on indices, not the actual $x$ positions.

Something very interesting has happened to our Fourier coefficients as a direct result of discretizing the data.  Notice that:
$$c_{k+n} = c_{k}$$
so evidently we have only $n$ unique Fourier coefficients.

Blah Blah Blah... Nyquist... Blah Blah blah.

\begin{eqnarray*}
  \psi(x) &=& \sum_{j=j_a}^{j_b} c_j e_j (x) \\
  &=& \frac{1}{\sqrt{a}}\sum_{j=j_a}^{j_b} c_j  \exp\left(i \frac{2\pi j x}{a}\right) \\
  &=& \frac{1}{\sqrt{n}}\sum_{j=j_a}^{j_b} A_j  \exp\left(i \frac{2\pi j (x-x_0)}{a}\right) \\  
\end{eqnarray*}


\section{Homework Problems}

\noindent
{\bf Problem 1:} Use python {\tt fft.fft} to calculate FS for a sine and cosine.  Plot the magnitude and phase at each wave number and compare to what you expect.


{\bf Problem 2:} Plot $f(x)$ by expanding using the Fourier coefficents from Problem 1.
Do not use inverse discrete fourier transform!  We want a smooth function!


{\bf Problem 3:} For square wave, plot coefficents, then plot f(x) obtained from expansion in fourier coefficents/

...
    


    

    

    




\noindent
{\bf Problem 2:} Write a program to determine the coefficients of the Hermite polynomial $H_n(u)$ for any degree $n$ by using the generating function:
$$\exp(-z^2+2zu) = \sum_{n=0}^{\infty} \frac{z_n}{n!}H_n(u)$$
Print the non-zero coefficients for $n=0$ to $5$ and
compare to e.g. Griffiths Table 2.1\\[5pt]

\noindent
{\bf Problem 3:} (A) Write a program to calculate the Hermite polynomial
$H_n(u)$ for a particular value $u$ and for any degree
$n$ by using the recursion formula:
$$H_{n+1}(u) = 2 u H_n(u) - 2 n H_{n-1}(u)$$.\\

\noindent
(B) Test your function {\bf thoroughly} against the Python library results
(scipy.special.hermite) using randomly chosen values.  That is, pick a
random integer $n$ from $0$ to some reasonable maximum test value $N$, then
pick a randomly chosen $u$ value from some reasonably range,
and compare the results of your program with that of the python library.  You
want to do this many many times, so you'll have to figure out a way to
handle the fact the agreement will not be perfect and yet still say
something definitive.\\[5pt]

\noindent
{\bf Problem 4:} Plot the wave function $\psi_n(x)$ and
$|\psi_n(x)|^2$ for $n=0,1,2,3$.  You are aiming to reproduce Fig
2.7(a) in Griffiths.  (Let me know if you don't have access to this textbook
and I will include photocopies of the needed figures on the course
site.)  Remember that $\psi_n(x)$ is not just the Hermite polynomical!
You may use the python library for the Hermite polynomials.\\[5pt]

\noindent
{\bf Problem 5:} Numerically verify that the $\psi_n(x)$ are
orthonormal.  Keep to $n<10$.  (Depending on your approach, you won't
get exactly 1 and 0 so you will need some way to handle this and still
give a definitive answer...there are many valid approaches)\\[5pt]

\end{document}
\documentclass[12pt]{book}

\usepackage[dvips,letterpaper,margin=0.75in,bottom=0.5in]{geometry}
\usepackage{cite}
\usepackage{slashed}
\usepackage{graphicx}
\usepackage{amsmath}
\usepackage{amssymb}
\usepackage{braket}
\begin{document}

\title{The Discrete Fourier Transform \\} 
\author{Michael Mulhearn}

\maketitle

\appendix
\setcounter{chapter}{6}
\chapter{The Discrete Fourier Transform}






\end{document}





